% GENERATED FILE. Edit canonical lessons, then run npm run generate:books.
% canonical-source-hash: fnv1a64:e64adf34317d3569
% canonical-lessons: RU-C385-obrazovanie, RU-C385-dissertaciya, RU-C385-kafedra, RU-C385-bakalavr, RU-C385-magistr

\chapter{Higher education}
\label{ch:ru-higher-education}
\hlchaptermodality{385}

\begin{chapteropening}
\textbf{By the end of this chapter:} \emph{I can talk about education and degrees: a bachelor and a master, a dissertation, and a university department.}

Complete the last lesson of chapter 385: I can talk about education and degrees: a bachelor and a master, a dissertation, and a university department.
\end{chapteropening}

\section[\texorpdfstring{\ru{образование} (obrazovániye)}{образование (obrazovániye)}]{\ru{образование} (obrazovániye) — education}

\label{lesson:RU-C385-obrazovanie}



\begin{quote}
Before the new one: say the Russian for a crib sheet, then the Russian for to get distracted.
\end{quote}

\subsection*{You'll want to know: \ru{образование}}
\textbf{\ru{образование}} (\emph{obrazovániye}) — \textquotedblleft{}education\textquotedblright{}.

\begin{quote}
\ru{У} \ru{неё} \ru{хорошее} \ru{образование}.
\end{quote}

\begin{quote}
\emph{U neyó khorósheye obrazovániye.}
\end{quote}

\begin{quote}
She has a good education.
\end{quote}

Neuter. \emph{\ru{Высшее} \ru{образование}} — higher education. It also means formation: \emph{\ru{образование} \ru{льда}}.

\begin{cousinweb}
\textbf{\ru{образование}} is built on \emph{\ru{образовать}} \textquotedblleft{}to form, to shape\textquotedblright{}, from \emph{\ru{образ}} \textquotedblleft{}a shape, an image\textquotedblright{}, inherited from Proto-Slavic \emph{*obrazŭ} (EtymDB). Education gives a mind its shape.
\end{cousinweb}

\begin{grammarlens}[title={What you've built}]
One more word for telling, explaining and describing at B1.
\end{grammarlens}

\subsection*{Your turn}
\begin{itemize}
\raggedright
  \item \emph{Say it:} \emph{obrazovániye}
  \item \emph{Say it:} \emph{obrazovániye}, once more
  \item \emph{Say it:} a sentence of your own with \emph{obrazovániye}
  \item \emph{Recall:} say the Russian for a crib sheet, then the Russian for to get distracted, then say \emph{obrazovániye} again
\end{itemize}

\begin{tcolorbox}[breakable,colback=teal!4,colframe=teal!35!black,title={Before you move on}]
What does \ru{образование} mean? (\textquotedblleft{}Education\textquotedblright{}.) Say it once more.
\end{tcolorbox}
\section[\texorpdfstring{\ru{диссертация} (dissertátsiya)}{диссертация (dissertátsiya)}]{\ru{диссертация} (dissertátsiya) — a dissertation, a thesis}

\label{lesson:RU-C385-dissertaciya}



\begin{quote}
Before the new one: say the Russian for to get distracted, then the Russian for education.
\end{quote}

\subsection*{You'll want to know: \ru{диссертация}}
\textbf{\ru{диссертация}} (\emph{dissertátsiya}) — \textquotedblleft{}a dissertation, a thesis\textquotedblright{}.

\begin{quote}
\ru{Он} \ru{пишет} \ru{диссертацию} \ru{о} \ru{русской} \ru{истории}.
\end{quote}

\begin{quote}
\emph{On píshet dissertátsiyu o rússkoy istórii.}
\end{quote}

\begin{quote}
He is writing a dissertation on Russian history.
\end{quote}

Feminine. \emph{\ru{Защитить} \ru{диссертацию}} — to defend a thesis.

\begin{cousinweb}
\textbf{\ru{диссертация}} is usually traced to Latin \emph{dissertātiō} \textquotedblleft{}a discourse\textquotedblright{} (Lewis \& Short), from \emph{dissertāre} \textquotedblleft{}to discuss\textquotedblright{}, the frequentative of \emph{disserere} \textquotedblleft{}to set out in order\textquotedblright{}, built on \emph{serere} \textquotedblleft{}to join, to string together\textquotedblright{}.
\end{cousinweb}

\begin{grammarlens}[title={What you've built}]
One more word for telling, explaining and describing at B1.
\end{grammarlens}

\subsection*{Your turn}
\begin{itemize}
\raggedright
  \item \emph{Say it:} \emph{dissertátsiya}
  \item \emph{Say it:} \emph{dissertátsiya}, once more
  \item \emph{Say it:} a sentence of your own with \emph{dissertátsiya}
  \item \emph{Recall:} say the Russian for to get distracted, then the Russian for education, then say \emph{dissertátsiya} again
\end{itemize}

\begin{tcolorbox}[breakable,colback=teal!4,colframe=teal!35!black,title={Before you move on}]
What does \ru{диссертация} mean? (\textquotedblleft{}A dissertation, a thesis\textquotedblright{}.) Say it once more.
\end{tcolorbox}
\section[\texorpdfstring{\ru{кафедра} (káfedra)}{кафедра (káfedra)}]{\ru{кафедра} (káfedra) — a university department; a lectern}

\label{lesson:RU-C385-kafedra}



\begin{quote}
Before the new one: say the Russian for education, then the Russian for a dissertation.
\end{quote}

\subsection*{You'll want to know: \ru{кафедра}}
\textbf{\ru{кафедра}} (\emph{káfedra}) — \textquotedblleft{}a university department; a lectern\textquotedblright{}.

\begin{quote}
\ru{Он} \ru{работает} \ru{на} \ru{кафедре} \ru{физики}.
\end{quote}

\begin{quote}
\emph{On rabótayet na káfedre fíziki.}
\end{quote}

\begin{quote}
He works in the physics department.
\end{quote}

Feminine. A department within a faculty: \emph{\ru{на} \ru{кафедре}}. It is also the lecturer’s desk.

\begin{cousinweb}
\textbf{\ru{кафедра}} is usually traced, through Latin \emph{cathedra}, to Greek \emph{kathédra} \textquotedblleft{}a seat, a chair\textquotedblright{}, built from \emph{katá} \textquotedblleft{}down\textquotedblright{} and \emph{hédra} \textquotedblleft{}a seat\textquotedblright{}: the teacher’s chair, then the department round it.
\end{cousinweb}

\begin{grammarlens}[title={What you've built}]
One more word for telling, explaining and describing at B1.
\end{grammarlens}

\subsection*{Your turn}
\begin{itemize}
\raggedright
  \item \emph{Say it:} \emph{káfedra}
  \item \emph{Say it:} \emph{káfedra}, once more
  \item \emph{Say it:} a sentence of your own with \emph{káfedra}
  \item \emph{Recall:} say the Russian for education, then the Russian for a dissertation, then say \emph{káfedra} again
\end{itemize}

\begin{tcolorbox}[breakable,colback=teal!4,colframe=teal!35!black,title={Before you move on}]
What does \ru{кафедра} mean? (\textquotedblleft{}A university department; a lectern\textquotedblright{}.) Say it once more.
\end{tcolorbox}
\section[\texorpdfstring{\ru{бакалавр} (bakalávr)}{бакалавр (bakalávr)}]{\ru{бакалавр} (bakalávr) — a bachelor (degree holder)}

\label{lesson:RU-C385-bakalavr}



\begin{quote}
Before the new one: say the Russian for a dissertation, then the Russian for a university department.
\end{quote}

\subsection*{You'll want to know: \ru{бакалавр}}
\textbf{\ru{бакалавр}} (\emph{bakalávr}) — \textquotedblleft{}a bachelor (degree holder)\textquotedblright{}.

\begin{quote}
\ru{После} \ru{четырёх} \ru{лет} \ru{учёбы} \ru{он} \ru{стал} \ru{бакалавром}.
\end{quote}

\begin{quote}
\emph{Pósle chetyryókh let uchyóby on stal bakalávrom.}
\end{quote}

\begin{quote}
After four years of study he became a bachelor.
\end{quote}

Masculine. \emph{\ru{Степень} \ru{бакалавра}} — a bachelor’s degree.

\begin{cousinweb}
\textbf{\ru{бакалавр}} is usually traced to Medieval Latin \emph{baccalaureus}, a reshaping of \emph{baccalārius} \textquotedblleft{}a young man, a junior\textquotedblright{}; the further origin is not settled.
\end{cousinweb}

\begin{grammarlens}[title={What you've built}]
One more word for telling, explaining and describing at B1.
\end{grammarlens}

\subsection*{Your turn}
\begin{itemize}
\raggedright
  \item \emph{Say it:} \emph{bakalávr}
  \item \emph{Say it:} \emph{bakalávr}, once more
  \item \emph{Say it:} a sentence of your own with \emph{bakalávr}
  \item \emph{Recall:} say the Russian for a dissertation, then the Russian for a university department, then say \emph{bakalávr} again
\end{itemize}

\begin{tcolorbox}[breakable,colback=teal!4,colframe=teal!35!black,title={Before you move on}]
What does \ru{бакалавр} mean? (\textquotedblleft{}A bachelor (degree holder)\textquotedblright{}.) Say it once more.
\end{tcolorbox}
\section[\texorpdfstring{\ru{магистр} (magístr)}{магистр (magístr)}]{\ru{магистр} (magístr) — a master’s graduate}

\label{lesson:RU-C385-magistr}



\begin{quote}
Before the new one: say the Russian for a university department, then the Russian for a bachelor.
\end{quote}

\subsection*{You'll want to know: \ru{магистр}}
\textbf{\ru{магистр}} (\emph{magístr}) — \textquotedblleft{}a master’s graduate\textquotedblright{}.

\begin{quote}
\ru{Она} \ru{хочет} \ru{получить} \ru{степень} \ru{магистра}.
\end{quote}

\begin{quote}
\emph{Oná khóchet poluchít stépen magístra.}
\end{quote}

\begin{quote}
She wants to get a master’s degree.
\end{quote}

Masculine. \emph{\ru{Магистратура}} — a master’s course.

\begin{cousinweb}
\textbf{\ru{магистр}} is usually traced to Latin \emph{magister} \textquotedblleft{}a master, a teacher, a chief\textquotedblright{} (Lewis \& Short).
\end{cousinweb}

\begin{grammarlens}[title={What you've built}]
One more word for telling, explaining and describing at B1.
\end{grammarlens}

\subsection*{Your turn}
\begin{itemize}
\raggedright
  \item \emph{Say it:} \emph{magístr}
  \item \emph{Say it:} \emph{magístr}, once more
  \item \emph{Say it:} a sentence of your own with \emph{magístr}
  \item \emph{Recall:} say the Russian for a university department, then the Russian for a bachelor, then say \emph{magístr} again
\end{itemize}

\begin{tcolorbox}[breakable,colback=teal!4,colframe=teal!35!black,title={Before you move on}]
What does \ru{магистр} mean? (\textquotedblleft{}A master’s graduate\textquotedblright{}.) Say it once more.
\end{tcolorbox}
